\subsection{Constructivo por Invariante de Paridad (permutación con signos)}

\graybold{Preguntas guía:}

\begin{itemize}
\item ¿El problema pide construir explícitamente una configuración (o probar que es imposible), en vez de optimizar o contar?
\item ¿Puedo idear una construcción simple (simetrías, intercambios, patrones repetidos) en vez de buscar entre todas las opciones?
\item ¿La imposibilidad de otros valores se puede demostrar con un invariante (paridad, módulo, suma fija) que ninguna elección puede romper?
\end{itemize}

\graybold{Utilidad:} Resolver problemas donde no hace falta ningún algoritmo de búsqueda: basta con proponer una construcción explícita y justificar, con un argumento de invariante, que el resultado contrario (o cualquier otro valor puntual) nunca puede ocurrir sin importar qué elecciones se hagan.

\graybold{Complejidad:} \bigO{n} por caso - construir y emitir la respuesta.

\graybold{Fundamento teórico:} Se particiona \{1,...,n\} en pares consecutivos (2k−1, 2k) y se intercambian entre sí. Cada par contribuye al contador un valor relacionado con V\_k = (2k−1)(2k), que siempre es par; y ese aporte entra a la suma final como −V\_k, 0 o V\_k (las tres opciones son pares). Como la suma de términos pares es par, el resultado nunca puede ser impar - en particular, nunca puede ser 1.

\graybold{Ejercicio aplicado}

(Codeforces 2246A — farmpiggie and Subset Sum) Dado un entero par n, construir una permutación p de longitud n tal que, sin importar qué signo (sumar, restar o ignorar) se le asigne a cada término i·p\_i al acumular un contador, el resultado final nunca pueda ser exactamente 1.

\graybold{I/O:} t ≤ 25, n ≤ 50 (par) → entrada pequeña, readline por línea (patrón 1) es suficiente.

\graybold{Código solución}

\begin{lstlisting}[language=Python]
import sys
input = sys.stdin.buffer.readline
 
t = int(input())
out = []
for _ in range(t):
    n = int(input())
    p = []
    for i in range(1, n + 1):
        # empareja (1,2), (3,4), ..., (n-1,n) e intercambia cada pareja:
        # el aporte de cada pareja al contador siempre es par (ver teoria)
        if i % 2 == 1:
            p.append(i + 1)      # posicion impar -> apunta a la siguiente
        else:
            p.append(i - 1)      # posicion par -> apunta a la anterior
    out.append(' '.join(map(str, p)))
sys.stdout.write('\n'.join(out) + '\n')
\end{lstlisting}
